\documentclass{article}
\usepackage{graphicx} % Required for inserting images
\usepackage{syntax}
\usepackage{listings}
\lstset{
  basicstyle=\ttfamily\small,
  columns=fullflexible,
  keepspaces=true
}

\title{Mangofn formalization}
\author{Mangofn Development Team}
\date{August 2025}

\begin{document}

\maketitle

\section{Preface}
\section{Introduction}
\section{Syntax of Mangofn}
\subsection{Reserved Words}
The following are \textit{reserved words} in Mangofn. These keywords are part of the language syntax and may not be used as identifiers.

\subsubsection*{Layout Constructs}
\begin{itemize}
    \item \texttt{window}
    \item \texttt{column}
    \item \texttt{row}
    \item \texttt{border}
\end{itemize}

\subsubsection*{UI Elements}
\begin{itemize}
    \item \texttt{button}
    \item \texttt{text}
    \item \texttt{textbox}
    \item \texttt{checkbox}
    \item \texttt{radiobutton}
    \item \texttt{toggleswitch}
    \item \texttt{calendar}
\end{itemize}

\subsubsection*{Boolean \& Control Values}
\begin{itemize}
    \item \texttt{true}, \texttt{false}
    \item \texttt{hidden}, \texttt{wrap}
\end{itemize}

\subsubsection*{Styling Keywords}
\begin{itemize}
    \item \texttt{fontfamily}, \texttt{fontsize}, \texttt{fontweight}, \texttt{fontstyle}
    \item \texttt{lineheight}, \texttt{textalign}, \texttt{textwrap}, \texttt{texttrim}
    \item \texttt{color}, \texttt{bgcolor}, \texttt{margin}, \texttt{width}, \texttt{height}
    \item \texttt{corner}, \texttt{density}, \texttt{fill}, \texttt{hug}
    \item \texttt{italic}, \texttt{underline}, \texttt{strikethrough}
    \item \texttt{center}, \texttt{left}, \texttt{right}
\end{itemize}

\subsubsection*{Text Behavior}
\begin{itemize}
    \item \texttt{word}, \texttt{character}, \texttt{notrim}
\end{itemize}

\subsubsection*{Color Literals}
\begin{itemize}
    \item \texttt{red}, \texttt{blue}, \texttt{green}, \texttt{yellow}, \texttt{pink}, \texttt{black}, \texttt{white}
\end{itemize}

\subsubsection*{Language Constructs}
\begin{itemize}
    \item \texttt{let}, \texttt{function}
\end{itemize}

\subsubsection*{Symbols}
\begin{itemize}
    \item \texttt{\{ \}}, \texttt{( )}, \texttt{:}, \texttt{,}, \texttt{=}
\end{itemize}

\section{Lexical Grammar}
\begin{lstlisting}[language={},caption={Lexical rules}]
<ID> ::=  [a-zA-Z] [a-zA-Z0-9_]* ;
<INT> ::= 0 |  [1-9] [0-9]* ;
<STRINGLIT> : '"' (~["\\] | '\\' .)* '"' ;
<BOOL> ::= 'true' | 'false' ;
<HEX> ::= '#' [0-9a-fA-F]{8} ;
\end{lstlisting}

\subsection{Lexical Conventions}
Mangofn source files are interpreted as sequences of ASCII characters. Whitespace and comments are ignored except where required to separate tokens. Identifiers and keywords are case-sensitive. 

\section{Syntactic Grammar}
\begin{verbatim}
Prog        = Window  [ "{" { UIElem } { Func } "}" ] ;
Window = "window" STRINGLIT [ INT INT [ STRINGLIT ] ] ;

UIElem      = "button" STRINGLIT [ "{" { Prop }+ "}" ]
             | "text" STRINGLIT [ "{" { Prop }+ "}" ]
             | "textbox" STRINGLIT
             | "checkbox" STRINGLIT
             | "radiobutton" STRINGLIT
             | "toggleswitch" STRINGLIT
             | "calendar"
             | "togglebutton"
             | "border" "{" { Prop }+ UIElem "}"
             | "column" "{" { Prop }+ { UIElem } "}"
             | "row" "{" { Prop }+ { UIElem } "}" ;

Prop        = PropName ":" PropValue
             | "onclick" ( ":" ID | "->" "{" { Stmt }+ "}" | "->" Stmt ) ;

PropName    = "bgcolor" | "wrap" | "corner" | "color" | "density" | "hidden"
             | "width" | "height" | "margin" | "id" | "fontfamily" | "fontsize"
             | "fontweight" | "fontstyle" | "lineheight" | "textalign" | "textwrap"
             | "texttrim" | "label" | "onclick" ;

PropValue   = Color | BOOL | Thickness | Size | STRINGLIT | INT
             | FontStyle | TextAlign | TextWrap | TextTrim ;

Func        = "function" ID "()" "{" { Stmt }+ "}" ;

Stmt        = "set" "(" PropName "," ID "," Exp ")"
             | "update" "(" ID "," "{" { Prop }+ "}" ")"
             | Exp ;

Exp         = INT | STRINGLIT | ID | BOOL | ID "()" ;

Thickness   = INT | INT "," INT | INT "," INT "," INT "," INT ;
Color       = "red" | "blue" | "green" | "yellow" | "pink"
             | "black" | "white" | HEX ;
Size        = INT | "fill" | "hug" ;
FontStyle   = "underline" | "italic" | "strikethrough" ;
TextAlign   = "center" | "left" | "right" ;
TextWrap    = "overflow" | "wrap" | "forcewrap" ;
TextTrim    = "word" | "character" | "notrim" ;
\end{verbatim}

\section{Semantics of Mangofn}
\subsection{Denotational Semantics}
Start with the clean, compositional part:
\begin{verbatim}
[ button(label, props) ] : Env \rightarrow Node
[ column(props, elems) ] : Env \rightarrow Tree
\end{verbatim}
Each syntactic form denotes a pure value in your domain of UI trees. Here you can play with
\begin{enumerate}
\item Fix-point definitions (recursive layouts)
\item Domain equations (UI = elem x Props x [UI])
\item Functional mappings ([Prog] : Env -$>$ UIState)
\end{enumerate}

\subsection{Constraint / Attribute Grammar Semantics - Layout}
Model layout as attribute propagation or constraint solving.
\begin{verbatim}
attribute inherited availableWidth : Int
attribute synthesized layoutBoxes : [Rect]

rule Column:
	distribute(children, availableWidth)
\end{verbatim}
Use an attribute grammar formalism or an explicit constraint system:
\begin{verbatim}
child2.top = child1.bottom + margin
width(child) \leq parent.width
\end{verbatim}
This lets you experiment with:
\begin{enumerate}
\item Constraint satisfaction semantics
\item Fix-point evaluation of attribute grammars
\item Declarative layout solvers (Cassowary-style)
\end{enumerate}

\subsection{Operational Semantics - Dynamic behaviour}
Model the runtime transitions for interaction and statements.
\begin{verbatim}
( UI, Env, click(id) ) \rightarrow ( UI', Env' )
\end{verbatim}
or for command execution:
\begin{verbatim}
( set(prop, id, v), \sigma ) \rightarrow \sigma[id.prop \mapsto v]
\end{verbatim}
Here you can explore:
\begin{enumerate}
\item Small-step vs. big-step semantics
\item Confluence and determinism proofs
\item Modeling reactivity as a labelled transition system (LTS)
\end{enumerate}
\subsection{Static Semantics - Type Rules}
Specify rules like:
\begin{verbatim}
\gamma \turnstile width : Int \lor {fill, hug}
\gamma \turnstile onclick : Function | Block
\end{verbatim}
You can practice:
\begin{enumerate}
\item Typing derivations and inference
\item Rule-based validation of props
\item Developing a kind of system for style/value categories
\end{enumerate}

\subsection{Axiomatic Semantics - Reasoning About Programs}
Once you have an operational semantics, you can add Hoare-style reasoning:
\begin{verbatim}
{ hasProp(id, color, old) } set(color, id, red) { hasProp(id, color, red) }
\end{verbatim}
or temporal logic:
% - $\box(button.clicked \to \diamond(dialog.open))$

That gives you experience writing and proving correctness properties over UI behaviour. 

\subsection{Abstract Interpretation - Static Analysis Playground}
Use this layer to practice static reasoning or optimization techniques:
- Prove that all color balues are valid hex codes.
- Infer possible sizes of elements before layout.
- Approximate event dependency graphs.

This connects formal semantics with program analysis.
%\subsection{Evaluation Model}
%\subsection{Event Handling}
%\subsection{Layout Resolution}

\appendix
\section{Examples}
\begin{verbatim}
window "Mango UI Test" 600 400 "icon.png" {
    text "Welcome to Mango UI!" {
        fontsize: 40
        color: #ff00ff
    }
    column {
        margin: 0

        button "Click me" {
            width: 1000
        }
        button "Another button"
        checkbox "Enable advanced mode"
    }
    row {
        margin: 20

        text "Horizontal layout preview" {
            color: white
            bgcolor: red
        }
        radiobutton "Option A"
        text "Additional details" {
            color: yellow
            bgcolor: blue
        }
    }
}

\end{verbatim}

\end{document}
